对协方差自适应后端来说，这种失配会直接变成求解困难。以 CMA-ES 为例，同一子问题若同时包含极平坦和极陡峭的方向，优化器就会遇到步长两难：足以沿平坦谷底推进的步长对陡峭方向过大，而对陡峭方向不致发散的步长又使平坦方向推进极慢。协方差自适应能够逐步缓解这一失配，但它需要充足的采样与学习；在有限预算下，分组决定的局部条件性仍可能实质性增加后端求解难度。相比之下，现有分组研究很少具体讨论优化器实际面对的这种局部条件性。